\documentclass[journal,10pt]{IEEEtran}

% The supplementary material for "Repo as Root Disk". Everything here is
% referenced from the main text by its S-number. Same class and packages as
% the paper, so the algorithms and tables set identically.
\usepackage{cmap}
\usepackage[T1]{fontenc}
\usepackage{amsmath}
\usepackage{booktabs}
% Paragraph columns set ragged right: justified narrow cells stretch the
% spaces between words into gaps.
\usepackage{array}
\newcolumntype{P}[1]{>{\raggedright\arraybackslash}p{#1}}
\usepackage{algorithm}
\usepackage{algpseudocode}
\usepackage{url}
\usepackage{cite}
\usepackage{tikz}
% The main text's palette, for the figure moved here from it.
\definecolor{ink}{HTML}{1A1A1A}
\definecolor{rule}{HTML}{9A9A9A}
\definecolor{accent}{HTML}{1F5C4D}
\usepackage[hidelinks]{hyperref}
\hypersetup{pdftitle={Supplementary Material for Repo as Root Disk: A Code Hosting Service as the Backing Store for Browser-Native Virtual Machines}, pdfauthor={Prasad Patil}}

% Pointers into the main text are generated from its label file by xref.mjs,
% never typed (AGENTS.md 3.3): \Mref{label} prints the main text's own number.
\makeatletter
\newcommand{\Mref}[1]{\@ifundefined{xr@M@#1}{\textbf{??}\@latex@warning{Reference `#1' in the main text undefined}}{\@nameuse{xr@M@#1}}}
\makeatother
\InputIfFileExists{xref-main.tex}{}{}

% Every experimental number is generated from traces/ by results.mjs, never
% typed (AGENTS.md 3.1): \R{key} prints it. An unknown key prints ?? and warns.
\makeatletter
\newcommand{\R}[1]{\@ifundefined{r@#1}{\textbf{??}\@latex@warning{Result `#1' undefined}}{\@nameuse{r@#1}}}
\makeatother
\InputIfFileExists{generated/results.tex}{}{}
% Generated table and figure bodies. \@@input is TeX's own input, which,
% unlike LaTeX's \input, is expandable and so may sit inside a tabular.
\makeatletter
\newcommand{\gen}[1]{\@@input generated/#1 }
\makeatother

\renewcommand{\thealgorithm}{S\arabic{algorithm}}
\renewcommand{\thetable}{S\arabic{table}}
\renewcommand{\thefigure}{S\arabic{figure}}
\renewcommand{\thesection}{S-\Roman{section}}

% A column heading on two lines, which fits a narrow column without
% shrinking the table's type.
\newcommand{\hd}[1]{\begin{tabular}[b]{@{}c@{}}#1\end{tabular}}


\begin{document}

\title{Supplementary Material for\\
Repo as Root Disk: A Code Hosting Service as the\\
Backing Store for Browser-Native Virtual Machines}
\author{Prasad~Patil}
\maketitle

\begin{abstract}
This supplement holds the material the main text refers to by S-number. It
gives the parameters of every experiment, the observed rate-limit envelope, the
restore and conflict-resolution algorithms, and the threat model. It also gives
the full contention grid with its live-service runs, a race of two guests on
one filesystem, the per-sync restore
tables for the captured histories, and the restic repeats. A command in the
artifact's README reproduces every number that has an archived trace; the next
section lists those that do not.
\end{abstract}

\section{Experiment Parameters}
\label{sup:params}

Table~\ref{tab:params} gives the parameters of every experiment in the main text.

\paragraph{What has no archived trace} The artifact's harnesses recompute most
figures in the paper from the traces it ships. The following were measured but
are not in an archived trace:
\begin{itemize}\itemsep1pt
\item the 2,460-entry tree GitHub's endpoint refused, and the 233,593-byte blob
GitLab's JSON endpoint returned once (archived reads did not see it again);
\item the 1\,GB machine's 299\,s commit to GitLab, whose record a later run
overwrote, and GitHub's throttling on the live replay beyond the per-sync rates
its log gives;
\item the other session's refusal at 541 a minute;
\item Codeberg's two 504 answers to a 256\,KB chunk, and the text of Forgejo's
403 ref-lock refusal, which an engine test keeps verbatim; the archived Forgejo
race records its outcome, not the message.
\end{itemize}

\paragraph{The loop in a browser} Headless Chrome with a fresh profile loads a
page that runs the loop of the main text's Section~\Mref{sec:loop} with the same modules, client
settings and guest as the Node harness. The driver speaks the DevTools protocol.
It counts every request the tab sends to the service from the browser's own
network events, where a preflight appears as an \texttt{OPTIONS} request of its
own, and assigns each to a phase by the wall-clock bounds the page records. The
token reaches the page as a function argument, never in a URL or a log. The
Node loop ran beside it in the same session. The two cache variants run the same
page with every request, or none, bypassing the browser's HTTP cache. GitHub
lets a reference be cached for \R{br.refMaxAge}\,s and updates it at another URL, so a tab left
to its cache rebuilt the state six syncs old in each of \R{br.none.rounds} rounds. Writes do not
bypass it, since that also skips the preflight cache: \R{br.all.preflights} preflights for \R{br.all.api}
requests.

\begin{table*}[!t]
\caption{Experiment parameters. In-process rows run the engine against a counting host, which isolates protocol cost from service variance; the others are measured against the live service. Every row runs under Node, with the emulator build the browser loads, except where marked: $^\dagger$ its capture ran in Chrome, $^\S$ it ran in Chrome, $^\ddagger$ it ran a Linux kernel under WSL2.}
\small
\centering
\setlength{\tabcolsep}{3pt}
\begin{tabular}{@{}lllP{0.94\columnwidth}@{}}
\toprule
experiment & disk & chunk & other parameters \\
\midrule
rate envelope & n/a & 4\,KB & 1 to 16 ways, 1\,s to 120\,s, \R{params.envWrites} writes \\
object latency & n/a & 16\,KB--4\,MB & 11 writes per size, interleaved, first round discarded \\
concurrency & n/a & 4, 256\,KB & 1 to 16 ways, four runs of 24 and of 16 writes \\
capture fidelity & 16\,MB & 256\,KB & in-memory and streamed buffers, \R{fid.perDevice} runs each \\
end-to-end loop & 16\,MB & 256\,KB & blank base, \R{loop.rounds} rounds on GitHub, \R{loop.increments} increments each \\
browser loop$^\S$ &16\,MB & 256\,KB & \R{br.browser}, \R{br.rounds} rounds beside \R{nb.rounds} under Node; cache variants \R{br.variantRounds} rounds each \\
read paths & n/a & 256\,KB & \R{reads.objects} objects, \R{reads.rounds} paired rounds, \R{reads.ways} ways \\
host comparison & 64\,MB & 256\,KB & 1 to 16 chunks, \R{hosts.rounds} rounds, back to back \\
deduplication & 256\,MB & 256\,KB & three states, \R{occ.machines} machines on each device, each compacted cold \\
distribution & 256\,MB & 256\,KB & Alpine 3.20.10, five packages, \R{alp.rounds} runs; chroot and root boot \\
commit shapes$^\dagger$ & 256\,MB & 256\,KB & 1,000 seeded and 43 captured syncs, in process \\
object store$^\dagger$ & 256\,MB & 256\,KB & 43 captured syncs over HTTP, loopback server \\
chunk sweep$^\dagger$ & 256\,MB & 32\,KB--1\,MB & \texttt{mke2fs} ranges; 43-phase session re-chunked \\
restic & 256\,MB & 256\,KB & 0.19.1, rest-server 0.14.0, compression off \\
native snapshot$^\dagger$ & 256\,MB & n/a & v86 save\_state, three rounds, all in Chrome \\
contention & 64\,MB & 256\,KB & 2 to 8 writers, 50 rounds, in process \\
contention, live & 16\,MB & 256\,KB & 2 to 8 writers, 3 rounds; GitHub, GitLab and Gitea, and Forgejo at 16\,KB chunks \\
two guests & \R{tg.diskMB}\,MB & 256\,KB & v86 guests, ext2, \R{tg.kinds} kinds of work, \R{tg.rounds} rounds, in-process host; \texttt{e2fsck} in WSL \\
1\,GB history$^\dagger$ & 1\,GB & 256\,KB & corpus workload, 43 syncs; GitHub and GitLab \\
kernel over NBD$^\ddagger$ & \R{kern.sizeMB}\,MB & 256\,KB & Linux \R{kern.release}, ext4, rebuilt by a second daemon; GitHub \\
preflight & n/a & n/a & OPTIONS to GitHub and GitLab, no credential \\
failures & \R{faults.diskMB}\,MB & 256\,KB & six phases, four kinds, plaintext and encrypted, in process \\
flush ablation & 16\,MB & 256\,KB & v86 guest, ext2, \R{flush.rounds} rounds, in-process host \\
flush ablation$^\ddagger$ & 64\,MB & 256\,KB & Linux \R{kern.release}, ext4, \R{flush.rounds} rounds, in-process host \\
encryption & 256\,MB & 256\,KB & both captured histories, 1\,GB captured again; plaintext and AES-GCM, in process \\
\bottomrule
\end{tabular}

\label{tab:params}
\end{table*}

\section{Rate-Limit Envelope}
\label{sup:envelope}

Table~\ref{tab:envelope} is the envelope of Section~\Mref{sec:tier} of the main text: short bursts from
rest pass at any concurrency, and whether a sustained rate is refused depends on
what preceded it. The second session was refused at \R{env2.refused} a minute after \R{env2.refusedSeconds}\,s, then,
after the minute's wait, passed \R{env2.later} for \R{env2.laterSeconds}\,s at the width at which the first was
refused at \R{env.refusedLow}. It began \R{env2.afterKernelMin} minutes after the kernel run and \R{env2.afterHeavyMin} after the other
GitHub runs, on the same account. A third session, unarchived, was refused at
541. Every archived row comes
from \texttt{envelope.mjs}. The first session's two runs are recorded in their
console logs, \texttt{traces/envelope-github.log} and
\texttt{envelope-github-ceiling.log}, since later runs overwrote the JSON file
each wrote; the second session is \texttt{envelope-github-session2.log} and
\texttt{.json}. The
rate is what the run achieved rather than what it asked for, and the duration is
how long it held before the deadline or the refusal.

\begin{table}[htbp]
\caption{The observed envelope. The first \R{env.rows1} rows are one session's two archived
runs, of \R{env.runA} and \R{env.runB} writes. The next \R{env2.rows} are a second archived session of
\R{env.runC} writes, in the order it ran: after the minute's wait that its refusal at
\R{env2.refused} imposed, \R{env2.later} a minute passed, the rate the first session was refused at.
The last row is a third session, in a browser page whose log was not kept.}
\small
\centering
\begin{tabular}{@{}lrrl@{}}
\toprule
rate & duration & concurrency & result \\
\midrule
\gen{tab-envelope.tex}
\midrule
\gen{tab-envelope2.tex}
\midrule
\textbf{541/min} & \textbf{40\,s} & \textbf{8} & \textbf{403} (third session) \\
\bottomrule
\end{tabular}
\label{tab:envelope}
\end{table}

Two documented limits bound writing, 80 a minute and 500 an
hour~\cite{ghratelimits}, and neither
fired: the first run sustained \R{env.heldRate} a minute for \R{env.heldMinutes} minutes and made \R{env.runA}
writes in all, in minutes. The ceiling the service's points system implies, five
points a write against 900 a minute, is 180, and it did not fire either. What
did was a refusal between \R{env.passRate} and \R{env.refusedLow} a minute in the first session and between
\R{env2.pass} and \R{env2.refused} in the second, and a refusal costs a minute. The second session's
pass at \R{env2.later} after its wait shows the limit is not a fixed rate: what a rate meets
depends on what preceded it.

Table~\ref{tab:latency} is the cost per write by payload, eleven writes at
each size, interleaved with the first round discarded. Every median carries a
bootstrap 95\% confidence interval from 2,000 resamples of the observations
themselves at a fixed seed. The samples are small and a request-metered
service's distribution is not normal, so an interval assuming normality would
understate it. A fixed seed means a published interval can be re-derived from
the archived observations rather than believed. The tail is reported
beside the median because a sync is many writes and pays the tail: at \R{lat.worstTailSize}
the p95 is \R{lat.worstTail} times the median, and every other size is inside \R{lat.otherTail}.

\textbf{The intervals are as much of the finding as the medians are.} At every
size the interval spans more than a quarter of the median, from \R{lat.spanLow} at \R{lat.spanLowSize}
to \R{lat.spanHigh} at \R{lat.spanHighSize}. Sizes eight or more times apart are separated in every case,
but three nearer pairs are not. The intervals for 16 and 64\,KB overlap by
\R{lat.overlap16v64}\,ms, and 1\,MB's, widened by the tail its p95 shows, overlaps both
2\,MB and 4\,MB. A single median over a handful of writes against a live
service is not a quantity anyone should design against. The trend by payload is
solid, and the order of neighboring sizes is not.

\begin{table}[!htb]
\caption{Cost of one object write by payload: medians of \R{lat.reps} with a
bootstrap 95\% interval, on GitHub, no concurrency.}
\small
\centering
\begin{tabular}{@{}lrrrrr@{}}
\toprule
payload & writes & median & 95\% CI & p95 & MB/s \\
\midrule
\gen{tab-latency.tex}
\bottomrule
\end{tabular}
\label{tab:latency}
\end{table}

Fig.~\Mref{fig:crossover} of the main text plots these rates.

We measured concurrency twice, four runs at each level: once with a 4\,KB
payload to isolate the request rate, and once with the 256\,KB chunk the engine
actually writes. At 4\,KB, concurrency 1, 2, 4, 8 and 16 gave \R{conc.rates}
writes a minute. Runs at one level differ by up to \R{conc.spread} of the
median (at \R{conc.spreadWidth} ways). The per-request median lies between \R{conc.medLow} and \R{conc.medHigh}\,ms,
and its 95\% intervals all fall inside \R{conc.ciLow} to \R{conc.ciHigh}\,ms. Those intervals pool the
requests of all four runs and treat them as independent, so they describe the
spread of requests rather than of sessions; the run range beside each rate is
the session-to-session figure. That last figure is what makes the scaling real
rather than an artifact of a warming connection: concurrency does the work, and
the server does not get slower. At 256\,KB the same levels gave \R{chunkconc.rates}
a minute, which is \R{chunkconc.mbs}\,MB/s.

\textbf{The two series part company at sixteen ways.} The small payload holds
its \R{conc.small16Median}\,ms median there; the chunk-sized one rises from \R{chunkconc.serialMedian} to
\R{chunkconc.sixteenMedian}\,ms, so its rate reaches \R{chunkconc.speedup} times serial against the small payload's
\R{conc.small16Speedup}. Requests still scale and bytes stop, which suggests that by sixteen ways
the link, not the service, is the constraint; we did not isolate it. The
256\,KB series is also the noisier of the two, spanning \R{chunkconc.eightLow} to \R{chunkconc.eightHigh} a minute at
eight ways against \R{conc.smallEightLow} to \R{conc.smallEightHigh} for the small one. A measurement that moves bytes
inherits the variance of whatever else is using the link. That is why the
figures above are medians of four runs, reported with their range.

\section{How Often Writers Meet}
\label{sup:collisions}

The cost of contention only matters at the rate contention happens. A sync's
duration $d$ is the window in which a second writer can arrive from the same
parent: \R{loop.firstS}\,s for the loop's commit in Section~\Mref{sec:loop} of the main text, \R{hosts.github.sec16}\,s for
sixteen chunks on GitHub. Two writers syncing every $T$
seconds collide with probability about $2d/T$, which at five-minute intervals
is \R{coll.loopPct} per sync for the first and \R{coll.bigPct} for the second. The
approximation is the union bound on two uniform arrivals and is loose above a
few percent. That is enough for the design decision it supports. With a
per-round cost growing as $N^2$ and a collision rate that rises with sync
duration, the concurrency model is handoff and forking rather than simultaneous
writers.

\section{Algorithms}
\label{sup:algorithms}

Algorithm~\ref{alg:restore} is the restore of Section~\Mref{sec:sync} of the main text; its loop bound is the
bounded-restore property as code. The cipher fits a manifest when both are
encrypted under the same salt or neither is. The MAC check applies to an
encrypted manifest, and a missing digest fails the digest check, except in a
manifest of the format before digests restored without a pin.
Algorithm~\ref{alg:conflict} is the conflict
resolution the evaluation's contention experiment measures; the final line's
budget is one by default and a parameter otherwise. A refused epoch is kept:
\textsc{Fork} starts a branch at $\mathit{head}_{\mathit{prev}}$ and the next
sync lands the epoch there. The stale set $S$ is cleared by hydrating the
chunks it names, or the whole machine.

Algorithm~\Mref{alg:sync} of the main text leaves two details to the host
adapter's \textsc{Commit}. Under encryption the manifest carries its MAC (Section~\ref{sup:deterministic}).
On a host whose commits apply changes to the tree they inherit, the same request
deletes the objects the next manifest drops, with the manifest in the last
commit of a split sync (Table~\ref{tab:batch-tree}). A refusal is a lost race only where it
means one: GitHub answers any validation failure with 422, so its 422 counts
only when the reference update returns it, while Gitea's 422 counts by its
message, \texttt{sha does not match}.

\begin{algorithm}[!htb]
\caption{Restore. Note the loop bound.}
\label{alg:restore}
\begin{algorithmic}[1]
\Require host $H$, branch $B$, cipher $E$, optional pinned digest $d$
\State $\mathit{ref} \gets \textsc{ResolveRef}(H, B)$
\State $T \gets \textsc{ReadTree}(H, \mathit{ref}.\mathit{tree})$
\State $m \gets \textsc{ReadObject}(H, T[\mathit{manifest}])$
\If{$d$ given \textbf{and} $\textsc{Sha256}(m) \neq d$}
  \State \textbf{abort} \Comment{substituted, or an older state}
\EndIf
\State $M \gets \textsc{Parse}(m)$
\If{$E$ does not fit $M$ \textbf{or} $\textsc{Mac}_{E}(M) \neq M.\mathit{mac}$}
  \State \textbf{abort} \Comment{no or wrong key; or $M$ altered}
\EndIf
\State $\mathit{disk} \gets M.\mathit{baseIsBlank}$ \textbf{?} $\mathbf{0}^{M.n}$ \textbf{:}
  $\textsc{Fetch}(M.\mathit{base})$
\For{$i \in \textsc{Indices}(M)$} \Comment{the live set; the history is never read}
  \State $c \gets \textsc{ReadObject}(H, M.\mathit{chunks}[i])$
  \Statex \hfill \Comment{an object two chunks share is read once}
  \If{$\textsc{Sha256}(c) \neq M.\mathit{digests}[i]$}
    \State \textbf{abort} \Comment{object ids are SHA-1; this is the check that matters}
  \EndIf
  \State $\mathit{disk}[\textsc{Extent}(i)] \gets E.\textsc{Decrypt}(c, i)$
  \Statex \hfill \Comment{bound to its position; see the threat model}
\EndFor
\State \Return $\mathit{disk}$
\end{algorithmic}
\end{algorithm}

\begin{algorithm}[!htb]
\caption{Conflict resolution, after a refused reference update.}
\label{alg:conflict}
\begin{algorithmic}[1]
\Require prepared set $P$, zeroed set $Z$, manifest $M_{\mathit{prev}}$ and
  $\mathit{head}_{\mathit{prev}}$ as at preparation, stale set $S$, rebase budget $b$
\State $\mathit{ours} \gets \{i : (i,\cdot) \in P\} \cup Z$
\State reload the branch, giving $M_{\mathit{theirs}}$ and a new $\mathit{head}$
\State $\mathit{theirs} \gets \{i : M_{\mathit{theirs}}.\mathit{chunks}[i] \neq
  M_{\mathit{prev}}.\mathit{chunks}[i]\}$ over $i$ in either manifest
  \Statex \hfill \Comment{a chunk they zeroed is absent, and changed}
\If{$\mathit{ours} \cap \mathit{theirs} \neq \emptyset$}
  \State $M \gets M_{\mathit{prev}}$; $\mathit{head} \gets \mathit{head}_{\mathit{prev}}$
  \Statex \hfill \Comment{so the branch refuses this epoch again until it is forked}
  \State \textbf{raise} \textsc{Conflict}$(\mathit{ours} \cap \mathit{theirs})$
  \Statex \hfill \Comment{two filesystem states, not two edits}
\EndIf
\State $S \gets S \cup \mathit{theirs}$ \Comment{in the manifest now, not on the device}
\State \Return commit again onto the new $\mathit{head}$ with budget $b - 1$
\end{algorithmic}
\end{algorithm}

\section{Threat Model}
\label{sup:threat}

Table~\ref{tab:threat} is the threat model Section~\Mref{sec:limits} of the main text discusses. The residual
column is the substance: each mitigation narrows an exposure rather than
removing it, and two of them depend on the service rather than on us.

\begin{table*}[!t]
\caption{Threat model.}
\small
\centering
\begin{tabular}{@{}P{0.12\textwidth}P{0.17\textwidth}P{0.32\textwidth}P{0.30\textwidth}@{}}
\toprule
asset & adversary & mitigation & residual \\
\midrule
disk contents & the service & client-side encryption under random nonces; free space is absent, not stored & metadata: sizes, timing, chunk counts, occupancy \\
\addlinespace
disk contents & another reader of a public repository & encryption & the chunk graph stays public \\
\addlinespace
chunk integrity & anyone able to write objects & SHA-256 digest per chunk, checked on read and required for every chunk of a current-format manifest; under encryption, each ciphertext bound to its chunk index and the manifest authenticated by an HMAC under a second key from the passphrase & a plaintext machine's manifest and tree remain SHA-1 unless its digest is pinned; a whole older manifest, authentic in its day, needs the pin \\
\addlinespace
credential & script or extension with DOM access & fine-grained token, one repository; kept in page memory, never stored & bearer token readable by any script on the page \\
\addlinespace
installed software & whoever supplies a package & same-origin delivery & no signature check offline \\
\addlinespace
deleted data & anyone with history access & compaction, then host GC & a secret written once persists until both \\
\bottomrule
\end{tabular}
\label{tab:threat}
\end{table*}

\section{Contention}
\label{sup:contention}

Table~\ref{tab:contention-full} is the complete grid behind
Section~\Mref{sec:contention} of the main text: writers racing to sync one
machine. Each row is 50 rounds of four chunks per writer on a 64\,MB disk, in
process against a fast-forward-only reference through the real engine. ``Hot''
puts a quarter of each writer's chunks in a four-chunk region all of them share.
A writer refused in one round would fork in use; the simulation abandons its
epoch before the next round, so rounds stay independent. The plans are seeded,
but the order in which writers reach the reference depends on asynchronous
hashing, so the tallies move by about a point between runs, as a race does.

\textbf{A fourth failure mode lies outside this grid.} The main text lists three
failure modes the grid exposes. A fourth lies in the resolution that takes the
branch's state after a refusal. A full hydrate must also write zeros where the
branch dropped a chunk. Otherwise a chunk the other writer deleted keeps this
device's bytes, the stale guard clears anyway, and the carried epoch re-commits
those bytes over the deletion. That is the overlapping write the refusal exists
to stop. A reload therefore marks every chunk the branch changed as stale,
deletions included, and hydration writes zeros where the branch has none; two
engine tests hold it.

\begin{table*}[t]
\caption{Writers racing to sync one machine from one parent. Outcomes are
shares of attempts; requests are per landed sync; wasted requests are those of
attempts that did not land.}
\footnotesize
\centering
\begin{tabular}{@{}rlrrrrrrr@{}}
\toprule
writers & sets & budget & first & rebased & refused & aborted & req./landed & wasted \\
\midrule
\gen{tab-contention.tex}
\bottomrule
\end{tabular}
\label{tab:contention-full}
\end{table*}

Table~\ref{tab:contention-live} is the same race on GitLab: three rounds per
row, two chunks per writer, every writer's traffic under one governor at 150
writes per minute. The probe reads the head manifest after every round and
checks that each landed writer's chunks are named in it. Three rounds were not:
both syncs of the first two two-writer rounds landed without a conflict, and one
writer's chunks were absent afterwards; the same at eight writers with a budget
of seven. GitLab's \texttt{last\_commit\_id} is validated before the commit is
made, so two syncs arriving together both pass and the second overwrites the
first. Table~\ref{tab:contention-github} is the same probe on GitHub. No round
lost an update, since the fast-forward-only reference update is atomic. A
budget of one rebase aborted exactly the losers of the second race, half at four
writers and three quarters at eight.

\textbf{Forgejo, raced.} The same table gives the probe on Forgejo, on
Codeberg, with \R{live.forgejo.chunkKB}\,KB chunks: Codeberg's gateway answered the first commit of a
256\,KB chunk with 504 twice, and what the race measures is the manifest's
refusal, not the chunks'. No round lost an update. One writer at a time,
Forgejo looks like GitLab: its contents API requires the blob a writer expects
to replace and refuses a writer already stale with 409, which is all the
one-at-a-time probe can see (archived with this run). Racing shows the
difference. Each change is committed in a scratch clone and pushed against the
head it started from. A writer that passed the check alongside another is
therefore refused one step later by git's own ref lock, 403 ``cannot lock
ref'', and nothing reaches the branch over a moved head. The engine classifies
that refusal as a lost race, so the loser rebases rather than failing; a test
holds it. At eight writers \R{live.forgejo.errors} attempts failed
outright with a bare 500: writes the service failed rather than lost, none
counted as landed, their epochs kept.

\textbf{And Gitea.} Forgejo was forked from Gitea and shares its contents API,
so the same probes ran on gitea.com, which is not one of the three services
compared elsewhere in the paper. One writer at a time Gitea refused a stale
writer as Forgejo does, but with 422 rather than 409, a third spelling of the
same refusal. Raced, at our own 256\,KB chunk, which gitea.com accepted where
Codeberg's gateway did not, no round lost an update, and at eight writers
\R{live.gitea.errors} attempts failed with a bare 500. Codeberg's 504 is therefore more likely a
limit of that deployment than of the API, though the two projects have
diverged enough that one run on the other cannot settle it.

\begin{table}[!htb]
\caption{The race on GitLab, three rounds per row, two chunks per writer.
Outcomes are shares of attempts; ``lost'' counts rounds in which a landed
writer's chunks were missing from the head manifest afterwards.}
\scriptsize
\setlength{\tabcolsep}{3pt}
\centering
\begin{tabular}{@{}rrrrrrrrr@{}}
\toprule
writers & budget & first & rebased & refused & aborted & \hd{req./\\landed} & \hd{s/\\round} & lost \\
\midrule
\gen{tab-contention-gitlab.tex}
\bottomrule
\end{tabular}
\label{tab:contention-live}
\end{table}

\begin{table}[!htb]
\caption{The race on GitHub, Forgejo and Gitea, three rounds per row, two
chunks per writer, Forgejo's at \R{live.forgejo.chunkKB}\,KB. No round on any of them lost an update;
``err.'' counts attempts the service failed with a server error.}
\scriptsize
\setlength{\tabcolsep}{3pt}
\centering
\begin{tabular}{@{}rrrrrrrrrr@{}}
\toprule
writers & budget & first & rebased & refused & aborted & err. & \hd{req./\\landed} & \hd{s/\\round} & lost \\
\midrule
\gen{tab-contention-github.tex}
\bottomrule
\end{tabular}
\label{tab:contention-github}
\end{table}

\section{Two Guests on One Filesystem}
\label{sup:twoguest}

This section asks whether the conflict rule keeps a real filesystem consistent
when two guests write to one machine at once. The races of
Section~\ref{sup:contention} commit synthetic chunk sets, which say nothing
about the filesystem the chunks hold.

\textbf{Setup.} One guest formatted a \R{tg.diskMB}\,MB disk with its own
\texttt{mke2fs}: ext2, 1\,KB blocks, eight block groups. It made two
directories of four 64\,KB files each and committed the disk at 256\,KB chunks.
For each kind of work in Table~\ref{tab:two-guest}, a branch was cut from that
commit and two guests booted from it, each under its own engine attached at
the branch head. They did the work on separate files, in separate kernels with
separate page caches. Guest A synced first and guest B second. A refused B
forked, as the application offers, and synced there. Capture was armed at the
prompt, as in the application, so each epoch held the write of the guest's
read-write mount. The host was the in-process one with fast-forward-only
references. \R{tg.rounds} rounds of \R{tg.kinds} kinds gave \R{tg.pairs} pairs
(\texttt{traces/two-guest.json}).

\textbf{Checks.} Every image was checked outside the guests, with
\texttt{e2fsck -fn} from e2fsprogs \R{tg.e2fsck} for the filesystem. Each file
a guest wrote was read back with \texttt{debugfs} and compared by MD5 with what
that guest had read after writing it, and the deleted file was looked up. Where
blocks had been allocated, the base image included, \texttt{e2fsck} remarked
that the superblock's free counts lagged the group descriptors. It recomputes
those counts and exits 0, and we count such an image as sound.

\begin{table*}[!t]
\caption{Two guests booted from one ext2 commit, each doing its half of the
work on a separate file, A syncing first; \R{tg.rounds} rounds a row. Chunks
are each epoch's; ``shared'' counts the chunks both changed, chunk 0 always
among them. ``Sound'' counts images of A's branch and B's fork that pass
\texttt{e2fsck -fn} with every file intact. The last column describes an image
the engine never commits: every chunk B changed laid over A's head, as a rebase
without the overlap check would commit.}
\small
\centering
\begin{tabular}{@{}lrrrrl@{}}
\toprule
work & \hd{chunks\\A\,/\,B} & shared & \hd{B\\refused} & sound & merge without the overlap check \\
\midrule
\gen{tab-two-guest.tex}
\bottomrule
\end{tabular}
\label{tab:two-guest}
\end{table*}

\textbf{Result.} Every second sync was refused, and chunk 0, which holds the
superblock and the group descriptors, was shared every time. A guest's sync
rewrites the superblock even after work that allocates nothing: the in-place
rewrites shared chunk 0 alone. All \R{tg.images} images of the branches and
forks were sound. On ext2, then, two mounted guests never take the rebase path,
which serves epochs that share no filesystem metadata.

\textbf{What the refusal prevents.} For each pair we also built two images the
engine never commits. The first lays every chunk B changed over A's head, as
Table~\ref{tab:two-guest} describes. The second keeps A's copy of each shared
chunk and merges only B's disjoint ones. It fared the same in every kind, losing
B's write where the first lost A's. Each guest allocated from the same bitmaps
and group descriptors, unaware of the other. A chunk-level merge therefore kept
one guest's allocation and dropped the other's, or kept both allocations of the
same blocks. Only the in-place rewrites, which allocate nothing, merged soundly,
and they were refused for the superblock alone. That refusal is conservative,
and a merge that understood the filesystem could admit it; we have not built one.

\textbf{Limits.} The experiment covers one filesystem, one kernel, one disk
geometry and small writes. A filesystem that does not rewrite shared metadata
at every sync could produce disjoint epochs, which would rebase, and this
experiment does not establish that the result is consistent. One active writer
per branch, with forks for parallel work, remains the recommendation.


\section{The Captured Histories}
\label{sup:histories}

Table~\ref{tab:history-full} gives every checkpoint of the captured \R{hist.syncs}-sync
history of Section~\Mref{sec:scaling} of the main text, on the 256\,MB machine:
format, Alpine, vim, then a forty-step working session.
Table~\ref{tab:history-1gb} gives the same for the 1\,GB machine under the
corpus workload. Table~\ref{tab:history-long} continues the 256\,MB session to
\R{long.syncs} syncs by repeating its working half: append, partial rewrite, vim edit,
delete and create, and package churn. Each covers the three commit shapes and restic, with the dirty
chunks of the sync and the live set at each point. The
engine represents a zeroed chunk on a blank base by absence and fetches each
distinct object once on restore, so the chunk-exploded column can be below the
live set where the guest copied files onto chunk-aligned blocks.

\begin{table}[!htb]
\caption{The \R{hist.syncs}-sync 256\,MB history: restore cost as requests / MB fetched.}
\scriptsize
\setlength{\tabcolsep}{2.5pt}
\centering
\begin{tabular}{@{}rrrrrrr@{}}
\toprule
sync & dirty & live & whole image & delta pack & chunk exploded & restic \\
\midrule
\gen{tab-history-256.tex}
\bottomrule
\end{tabular}
\label{tab:history-full}
\end{table}

\begin{table}[!htb]
\caption{The 1\,GB history under the corpus workload: restore cost as requests
/ MB fetched. At the last sync the stores hold \R{gb.whole.storageMB}, \R{gb.delta.storageMB}, \R{gb.ours.storageMB} and \R{gb.restic.storageMB}\,MB
and have spent \R{gb.whole.uploads}, \R{gb.delta.uploads}, \R{gb.ours.uploads} and \R{gb.restic.uploads} upload requests. On the services, the
machine committed to GitHub in \R{gb.github.commitRequests} requests and \R{gb.github.commitSeconds}\,s at 150 writes per
minute, the tree built in batches, and to GitLab in \R{gb.gitlab.commits} batch commits and
299\,s. It restored from GitHub in \R{gb.github.restoreRequests} requests: \R{gb.github.serialA} and \R{gb.github.serialB}\,s serially
and \R{gb.github.wide}\,s eight-wide, the run then meeting the hourly ceiling. It restored from
GitLab in \R{gb.gitlab.restoreRequests} requests (\R{gb.gitlab.treePages} of them tree pages, directories included),
\R{gb.gitlab.serial}\,s serially and \R{gb.gitlab.wide}\,s eight-wide, medians of three.}
\scriptsize
\setlength{\tabcolsep}{2.5pt}
\centering
\begin{tabular}{@{}rrrrrrr@{}}
\toprule
sync & dirty & live & whole image & delta pack & chunk exploded & restic \\
\midrule
\gen{tab-history-1gb.tex}
\bottomrule
\end{tabular}
\label{tab:history-1gb}
\end{table}

\begin{table}[!htb]
\caption{The 256\,MB session continued to \R{long.syncs} syncs: restore cost as requests
/ MB fetched, with restic measured every fiftieth sync. Delta pack's restore
first costs more requests than ours at sync \R{long.crossover}, \R{long.crossDelta} to \R{long.crossOurs}, and stays above it.
At the last sync the stores hold \R{long.whole.storageMB}, \R{long.delta.storageMB}, \R{long.ours.storageMB} and \R{long.restic.storageMB}\,MB and have
spent \R{long.whole.uploads}, \R{long.delta.uploads}, \R{long.ours.uploads} and \R{long.restic.uploads} upload requests.}
\scriptsize
\setlength{\tabcolsep}{2.5pt}
\centering
\begin{tabular}{@{}rrrrrrr@{}}
\toprule
sync & dirty & live & whole image & delta pack & chunk exploded & restic \\
\midrule
\gen{tab-history-long.tex}
\bottomrule
\end{tabular}
\label{tab:history-long}
\end{table}

Table~\ref{tab:batch-tree} replays the long session on an in-process host that
commits as GitLab does: each commit applies creations, updates and deletions to
the tree it inherits, and a listing costs a request per 100 entries,
directories included. The engine deletes from the tree what the next manifest does
not name. The variant beside it never deletes, so its tree keeps every
object ever written. The 1\,GB machine's GitLab restore, of a branch committed in
one pass, listed \R{gb.gitlab.treePages} pages, as this model with deletions does for it.

\begin{table}[!htb]
\caption{The long session on a batch-commit host: tree files, tree pages and
restore requests, with the engine's deletions and without them.}
\scriptsize
\setlength{\tabcolsep}{2.5pt}
\centering
\begin{tabular}{@{}rrrrrrrr@{}}
\toprule
& & \multicolumn{3}{c}{with deletions} & \multicolumn{3}{c}{without} \\
\cmidrule(lr){3-5}\cmidrule(lr){6-8}
sync & live objects & files & pages & restore & files & pages & restore \\
\midrule
\gen{tab-batch-tree.tex}
\bottomrule
\end{tabular}
\label{tab:batch-tree}
\end{table}

Table~\ref{tab:checkpoint} runs a delta pack that writes one pack of every chunk
holding data each \R{ck.every} syncs, and drops the packs before it from its tree, on
the seeded plan and the long session. Restore is measured just before and just
after a checkpoint, its worst and best cases.

\begin{table}[!htb]
\caption{Delta pack with a checkpoint every \R{ck.every} syncs: restore as requests / MB
fetched, storage held, and cumulative upload requests. Ours holds
\R{scale.ours.storageMB}\,MB at sync \R{scale.syncs} of the seeded plan, after \R{scale.ours.upload} upload requests.}
\scriptsize
\setlength{\tabcolsep}{3pt}
\centering
\begin{tabular}{@{}lrrrr@{}}
\toprule
workload & sync & restore & storage, MB & uploads \\
\midrule
\gen{tab-checkpoint.tex}
\bottomrule
\end{tabular}
\label{tab:checkpoint}
\end{table}

\textbf{The 1\,GB machine synced live.} Synced to GitLab one phase at a time,
the \R{gb.syncs} commits cost \R{gb.gitlab.liveRequests} requests and \R{gb.gitlab.liveMinutes} minutes, and the history then restored
in \R{gb.gitlab.liveRestore}\,s, byte-identical. On GitHub the first \R{gb.github.liveSyncs} commits cost \R{gb.github.liveRequests}
requests. The stream had by then spent some 2,200 requests in the hour, and the
service throttled it from 150 to between \R{gb.github.throttleLow} and \R{gb.github.throttleHigh} writes a minute, each
refusal costing its retry-after.

\section{Other Repository Operations}
\label{sup:operations}

\textbf{Compaction's read and GitLab's body size.} Compaction reads the disk back
first. On the streamed device, reading and hashing all 1,024 chunks took
\R{compact.streamedEmpty}\,s nearly empty and \R{compact.streamedFull}\,s filled, medians of \R{compact.rounds} cold runs. Written chunks
are local, while unwritten ones come from the base, here a local file and in a
browser a ranged request. On the in-memory blank it took \R{compact.memoryEmpty} and \R{compact.memoryFull}\,s.
GitLab accepted a \R{hosts.gitlab.oneRequestMB}\,MB request body, \R{hosts.gitlab.oneRequestPayloadMB}\,MB of chunks in base64, as one
commit. Our adapter split the next step, \R{hosts.gitlab.splitMB}\,MB, into \R{hosts.gitlab.splitRequests} commits at a ceiling we
chose rather than one the service published, so the service's limit is
unmeasured and our ceiling may split commits that need not be split. In the run
of Table~\Mref{tab:hosts} our governor halved its width \R{ours.halvings} times and never restored it,
and a variant that cannot restore it took \R{ratchet.sec16}\,s against \R{hosts.github.sec16}\,s, so that run
cannot show what recovery is worth.

A repository's other operations only partly survive the change of unit. Naming
a state does: a tag is a reference that never moves, one request and no
storage. Searching does, because every state costs the live set plus three to
restore whatever its age, so binary search over a history is affordable.
Finding which sync broke an environment across 64 commits took eight
reconstructions, against 64 for a scan. A replaying store, paying more for older
states, cannot do this, and neither can a versioned recipe, which has nothing to
boot. A fork hands someone a configured machine rather than sharing one. A pull
request is a proposed reference update with a gate, and the merge does not
survive, since overlapping chunk sets are refused. Diff and blame do not map at
block granularity.

Two behaviors of a batch-commit host appear only when it is synced twice. A
batch commit inherits the previous tree, so a commit that lists every live
object asks the host to create files that already exist. The adapter sends only
the files the tree lacks, and deletes those the next manifest does not name
(Table~\ref{tab:batch-tree}). And GitLab's \texttt{start\_sha} is not a parent pin: it means ``create
this branch from that commit'' and is refused for a branch that exists.

Committing the 1\,GB machine of Section~\Mref{sec:scaling} of the main text showed three
more. GitHub's tree endpoint fails at 2,460 entries where 941 succeed, so the
adapter builds trees incrementally through \texttt{base\_tree}. GitLab lists a
tree 100 entries a page, directories included, so the reference and tree that
cost GitHub three requests cost it \R{gb.gitlab.refAndTree}. And GitLab's JSON blob endpoint once
returned 233,593 bytes of a 262,144-byte chunk whose object was intact. The
per-chunk digest of Section~\Mref{sec:limits} of the main text, designed against
substitution, refused that restore, and the raw endpoint serves the bytes whole.
The truncation did not recur in \R{reads.gitlab.reads} later reads of random 256\,KB chunks
through the same endpoint.
% And a repository is not a shared drive: two writers mounting one
%device read-write corrupt it, as any block device without a cluster filesystem
%would, and that failure is silent.

\section{Chunk Size}
\label{sup:chunk}

The prototype fixes the chunk at 256\,KB. Table~\ref{tab:chunk} is what other
sizes would have cost on a captured trace of \texttt{mke2fs} on a 256\,MB disk,
\R{chunk.ranges} intercepted writes of about \R{chunk.writeKB}\,KB each at 8\,MB intervals, which is a
filesystem laying down block-group metadata. Nothing about the guest depends on
the chunk size: it writes the same bytes to the same offsets whatever the engine
later does with them. The dirty set is a pure function of the trace, so every
size is computed exactly from one run rather than measured across five. One run
per workload makes the sweep a screening, not a basis for choosing a size.

Two costs move in opposite directions. Requests fall as the chunk grows, since
a fixed set of writes lands in fewer chunks. Bytes rise, because a write that
straddles a boundary sends both chunks whole: a \R{chunk.writeKB}\,KB write costs two 256\,KB
chunks, which is the \R{chunk.amp256}$\times$ at our own size. The time column resolves
them against the rate the client paces to, 150 writes a minute, and the
\R{chunk.eightWay}\,MB/s measured eight-wide in Section~\ref{sup:envelope}, and shows where the binding
constraint changes hands.

\begin{table}[!htb]
\caption{What each chunk size would have cost on the \texttt{mke2fs} capture.
Time is the larger of the two costs at the 150 writes/minute the client paces
to and the \R{chunk.eightWay}\,MB/s measured eight-wide, and ``bound by'' says which.
Restore requests are arithmetic for the \R{chunk.liveMB}\,MB live set of Section~\ref{sup:histories}'s
256\,MB machine, $\lceil L/k \rceil + 3$.}
\small
\centering
\setlength{\tabcolsep}{2.2pt}
\begin{tabular}{@{}lrrrrrl@{}}
\toprule
chunk & dirty & sent & amplif. & \hd{write\\req} & \hd{restore\\req} & time \\
\midrule
\gen{tab-chunk.tex}
\bottomrule
\end{tabular}
\label{tab:chunk}
\end{table}

\textbf{On this trace our design point is not the optimum.} The write path is
request-bound up to 512\,KB and byte-bound at 1\,MB, so the cost is least at
512\,KB, where 256\,KB costs \R{chunk.formatRatio} times as long; and restore requests halve with
every doubling, which favors larger chunks again.

\paragraph{A second workload, and a different shape of answer}
That trace is a filesystem format, the workload most favorable to large
chunks: every write in it is a quarter-megabyte of bulk metadata. A working
session is the opposite, small scattered writes whose amplification rises with
the chunk, so we repeated the sweep on one. The \R{chunk2.phases}-phase session of
Section~\ref{sup:histories} was captured at 256\,KB, but a capture keeps more
than it was taken for. It stores every dirty chunk's contents, so the disk after
each phase can be reconstructed exactly. Two consecutive reconstructions then
differ in precisely the bytes that phase wrote. That recovers a byte-level trace,
compared a 4\,KB block at a time since a block is what a disk writes, and
every chunk size is then computed from it exactly by the same method. A write
that lays down the bytes already present leaves no difference and so is
invisible here; it is also free, because the engine hashes the chunk, finds it
present and sends nothing. A changed chunk is not always an upload either: its
new content may already be stored, from an earlier sync or another index, and
the engine's dedup probe then references it. Each size keeps the contents it
has stored, as the probe does over a session. As a check, at the capture's own
256\,KB the sweep gives \R{chunk2.uploads} uploads, and an independent count over the captured
chunks reproduces that exactly. The format trace of Table~\ref{tab:chunk} has ranges and no
contents, so it cannot model the probe and its sent column is an upper bound.

\begin{table}[!htb]
\caption{The same sweep on the \R{chunk2.phases}-phase working session, summed over the whole
run. Each phase is a sync and pays its own manifest, tree and commit; a chunk
ending a phase entirely zero is a deletion, and one whose content is already
stored a reference, rather than an upload. Limits as in Table~\ref{tab:chunk}.}
\small
\centering
\setlength{\tabcolsep}{3pt}
\begin{tabular}{@{}lrrrrl@{}}
\toprule
chunk & dirty & sent & amplif. & \hd{write\\req} & time \\
\midrule
\gen{tab-chunk2.tex}
\bottomrule
\end{tabular}
\label{tab:chunk2}
\end{table}

\textbf{The prediction holds and the conclusion survives it.} Amplification on
the working session is far worse at every size, \R{chunk2.amp256}$\times$ against
\R{chunk.amp256}$\times$ at our own 256\,KB and \R{chunk2.amp1M}$\times$ against \R{chunk.amp1M}$\times$ at
1\,MB, which is the scattered-write penalty the format trace could not show.
The crossover does not move: both workloads are request-bound through 512\,KB
and byte-bound at 1\,MB, so both put the least cost at 512\,KB, and on both
1\,MB is dearer than 512\,KB and cheaper than our 256\,KB. Our design point
costs \R{chunk.formatRatio} times the optimum on the format and \R{chunk2.ratio} times on the session.

We report this rather than tune to it. The analysis below 512\,KB is the same
analysis that justifies the design's central trade, which is to spend requests
rather than let restore grow, so a reader can see the trade rather than take
the constant on faith. A deployment that knew its workload should raise the
size; the engine takes it from the manifest and nothing in the protocol
depends on 256\,KB.

%Two effects this table does not carry. Free space is represented by absence, so
%a larger chunk makes a sparsely written machine store more: a chunk is occupied
%if any byte in it is. And a larger chunk coarsens the conflict unit, since two
%writers collide when they touch one chunk rather than one block.

\begin{table}[!htb]
\caption{Committing a machine that carries a distribution, medians of
\R{alp.rounds} runs. Installed is \texttt{apk}'s own count. Chunks are those holding
data; stored is their distinct objects at 256\,KB, which is what committing the
machine from blank uploads. Stored exceeds installed because a chunk is sent
whole however little of it changed.}
\label{tab:distribution}
\centering\small
\setlength{\tabcolsep}{3.5pt}
\begin{tabular}{@{}lrrrrr@{}}
\toprule
State & Installed & Chunks & Objects & Stored & Ratio \\
\midrule
\gen{tab-distribution.tex}
\bottomrule
\end{tabular}
\end{table}

\section{Staging Software Into the Guest}
\label{sup:staging}

The machine has no network by design, so software reaches it by a second path.
v86 exposes a 9p filesystem the guest mounts and the page writes files into
directly, so a rootfs archive or a package file travels from the page's own
origin into the guest without a request to any third party. That share lives in
browser memory and is empty after a reload, which makes it staging rather than
storage: anything worth keeping is copied onto the block device, where the sync
protocol can see it. Table~\ref{tab:distribution} is what committing such a machine costs.
The image's own guest is a Buildroot system built against uClibc, with no
package manager and no repository built against its C library, so the
distribution goes on the block device. A chroot into an Alpine minirootfs there
is a complete musl userland with a working \texttt{apk}, and every file of it
lands on the device the sync protocol commits.

\textbf{The machine of Section~\Mref{sec:alpine} of the main text.} We unpack an Alpine 3.20.10
minirootfs, \R{alp.archiveMB}\,MB across \R{alp.entries} entries, onto the device through the 9p share and
chroot into it; \texttt{apk-tools \R{alp.apk}} reports \R{alp.before} packages installed. Installing
\texttt{vim} from \R{alp.files} local \texttt{.apk} files totaling \R{alp.filesMB}\,MB, in one invocation,
reports \texttt{\R{alp.okLine}}, and \texttt{vim --version} runs inside the chroot.

Two details belong to the offline setting. Packages install as one selection,
since \texttt{apk} resolves only the set it is handed: \texttt{vim} alone fails
on \texttt{libncursesw} and succeeds with its dependencies named alongside. And
installation is unverified, since \texttt{--allow-untrusted} is what makes a
local file installable at all: a real reduction in assurance.

The \texttt{chroot} is a convenience, not a requirement. The same commit,
restored onto a disk, booted as the root filesystem of the image's own kernel,
\R{root.kernel}, started with \texttt{root=/dev/sda} and an \texttt{rdinit} naming no program, so
that the kernel mounts the root itself. The distribution's BusyBox \R{root.busybox} shell ran
as process~1 on the disk mounted at \texttt{/}, where \texttt{apk} listed \R{root.packages} packages,
\texttt{vim} ran and its file was present, in each of \R{alp.rounds} runs.

Committing the minirootfs stores \R{alp.rootfsObjects} objects, \R{alp.rootfsMB}\,MB, for the
\R{alp.rootfsInstalled}\,MB that \texttt{apk} counts as installed. Once \texttt{vim} brings that to
\R{alp.vimInstalledExact}\,MB, which its summary line rounds down to \R{alp.okMiB}\,MiB, the machine is \R{alp.vimObjects} objects and \R{alp.vimMB}\,MB.
The ratio of stored to installed bytes \emph{falls} as the machine grows, from
\R{alp.ratioRootfs}$\times$ to \R{alp.ratioVim}$\times$, the opposite of the expected direction; we have not
isolated why.

\section{Capture Fidelity}
\label{sup:fidelity}

The guest runs in the same emulator build under Node, so the checks are a
harness with an archived trace. Booting made \R{fid.bootWrites} writes to the disk, and one
\texttt{dd} of one sector arrived as one intercepted call of \R{fid.sectorKB}\,KB. Formatting a filesystem, mounting it, writing a
file and unmounting then produced \R{fid.calls} intercepted calls, \R{fid.kb}\,KB or about
\R{fid.kbPerCall}\,KB each, across stages that repeatedly rewrite the same structures.
The commit of that epoch, restored onto zeros, reproduced the live disk with
\emph{\R{fid.differing} differing bytes} across 16\,MB, so last-writer-wins ordering is
correct. To rule out luck, a second machine booted from nothing but the commit
mounted the filesystem and read back the file. On the streamed device, where no
single buffer exists to compare against, the device's own record of written
\R{fid.blockBytes}-byte blocks was the reference: \emph{\R{fid.blocks} of \R{fid.blocks}} reproduced.
That is the stronger check, since the reference comes from the device rather
than our interceptor. The write pattern was identical across all \R{fid.runs} runs,
\R{fid.perDevice} on each device.

The flush that precedes every seal was tested by leaving it out
(Section~\Mref{sec:kernel} of the main text). In each round the guest writes a
file holding a fresh marker, and two syncs follow at once, the first without
the flush and the second with it. The restored disk is then searched for the
marker. The same test runs on a Linux kernel through the NBD device, the
server in a container sharing the kernel and the host in process. Each round
also writes a second marker and polls the device every half second, unasked,
until it arrives. \texttt{flush-guest.mjs} and \texttt{flush-kernel.sh} archive both.

Capture also costs the guest nothing we could measure. In Chrome, on the machine
the captured history ends at, \R{wrap.mb}\,MB through the device took \R{wrap.attached}\,s with the
wrapper attached and \R{wrap.detached}\,s detached, medians of \R{wrap.pairs} interleaved rounds
spanning \R{wrap.low} to \R{wrap.high}\,s. A paired Wilcoxon signed-rank test over the rounds
finds no difference ($W = \R{wrap.W}$, two-sided $p = \R{wrap.p}$), though \R{wrap.pairs} rounds cannot
exclude a small cost.

\section{Further Limitations of the Prototype}
\label{sup:limits}

Four limits of the prototype matter to anyone deploying it. None of them
changes a measured result.

\paragraph{Installing offline gives up verification}
Removing the network keeps the no-server claim intact and moves the trust
decision to whoever put the package in front of the user. Serving packages from
the page's own origin makes that origin the trust anchor, a coherent position but
a weaker one than a signed index. A deployment that cared could verify signatures
in the page against Alpine's published keys before injection; we have not built
that.

\paragraph{The kernel is not versioned}
\label{sup:chroot}
The install runs in a chroot on the committed device, under a kernel and an
init from the read-only image. Section~\Mref{sec:alpine} of the main text also boots the
image's kernel with the committed device as its root and the distribution's
shell as process~1, which removes the image's userland. The kernel and its
command line still come from the image, so a machine cannot change its kernel,
and a kernel upgrade inside the distribution has no effect.

\paragraph{Machines have to be sized in advance}
The device is a fixed-size image, so a machine is provisioned rather than grown:
a 16\,MB disk holds the minirootfs and nothing useful after it, and the package
install above needs 256\,MB. Growing a machine is bookkeeping by design: the
manifest lists only written chunks, so its size can be raised, the guest
rebooted and the filesystem grown (Section~\Mref{sec:design} of the main text).
The prototype does not implement it, and refuses to attach a device whose size
differs from the manifest's; we have not measured a growth. Under-provisioning surfaces as an out-of-space error, not a
refusal.

\paragraph{What the budget permits}
The rate limit follows the engine onto the kernel: what must fit is the dirty
region between two syncs, about 187\,MB at 150 writes per minute, 256\,KB chunks
and a five-minute interval. Comfortable for a workstation, hopeless for a
database or a build server.

\section{Provider Capabilities}
\label{sup:capabilities}

The protocol's concurrency guarantee is not a property of the protocol. It is
the property of whichever atomic primitive the backing service exposes, and the
three services expose three different ones. Table~\ref{tab:capabilities} states
what each was measured to provide, and the guarantee that follows.

Git's own transport is not among them for a browser. We sent \R{gt.probes} probes, without
credentials, to the git endpoints of GitHub and GitLab: a preflight for a push,
and a cross-origin read of the push and fetch reference advertisements. None
carried \texttt{Access-Control-Allow-Origin}, and GitHub refused the preflight
outright, so a page can neither push nor read the answer
(\texttt{traces/git-transport.json}). Git's transport would otherwise push a
sync as one packfile and move the reference only from the value it expects, on
every host. A client outside the browser has no such limit, and for the kernel
device of Section~\Mref{sec:kernel} of the main text a push may be the cheaper
path; we have not measured it.

\begin{table*}[t]
\centering
\caption{Concurrency and atomicity capabilities of the evaluated services. The
guarantee is per commit: no landed sync is overwritten. Two guests sharing one
mounted filesystem are kept consistent by the refusal of the second, not by a
merge (Section~\ref{sup:twoguest}).}
\label{tab:capabilities}
\begin{tabular}{@{}P{0.11\textwidth}
                P{0.14\textwidth}
                P{0.18\textwidth}
                P{0.14\textwidth}
                P{0.09\textwidth}
                P{0.22\textwidth}@{}}
\toprule
\textbf{Service} &
\textbf{Object writes} &
\textbf{Atomic reference CAS} &
\textbf{Per-file CAS} &
\textbf{Lost updates} &
\textbf{Guarantee for simultaneous writers} \\
\midrule

GitHub &
Blob, tree, commit &
Yes: fast-forward-only reference update; stale parent refused (422) &
Not needed &
\R{live.github.lostOf} &
Holds: a lost race is detected; disjoint epochs rebase, overlapping ones are refused \\

\addlinespace

GitLab &
None; batch commit &
No: the commits API takes no expected parent for an existing branch, and a stale one is accepted &
Partial: \texttt{last\_commit\_id} on the manifest is refused when stale, but validated before the commit is made &
\R{live.gitlab.lostOf} &
\textbf{Does not hold}: two syncs arriving together both pass the check and the second overwrites the first \\

\addlinespace

Forgejo &
None; batch commit &
Not exposed, but enforced: the contents API pushes each commit against the head it started from &
Yes, required on every update: 409 to a writer already stale; a writer that goes stale mid-commit is refused at git's ref lock (403) &
\R{live.forgejo.lostOf} &
Holds, measured with up to eight writers; at eight, \R{live.forgejo.errors} attempts failed with a server error, none lost \\

\addlinespace

Gitea (the code Forgejo was forked from) &
None; batch commit &
Not exposed; its refusal of a writer that goes stale mid-commit was not observed separately &
Yes, required on every update: 422 to a writer already stale &
\R{live.gitea.lostOf} &
Holds, measured with up to eight writers at 256\,KB; at eight, \R{live.gitea.errors} attempts failed with a server error, none lost \\

\addlinespace

Object store &
Per-object PUT &
N/A &
Conditional write since 2024~\cite{s3conditional,s3cas2024}; S3 documents that of concurrent conditional writes the first to finish lands and the rest are refused &
Not measured on S3 &
Holds by S3's documented semantics, on the pointer; our loopback store refused the second of two writers (Section~\ref{sup:objectstore}) \\

\bottomrule
\end{tabular}
\end{table*}

The distinction that matters is between a compare-and-swap and a check-then-act.
GitHub's reference update compares and swaps in one operation, so a writer whose
parent has moved is refused. GitLab's \texttt{last\_commit\_id} compares first
and commits second. It refuses a writer that was already stale when it arrived
and cannot refuse one that goes stale in between. That is a narrower guarantee,
and the difference is the three lost rounds. Forgejo's is the atomic kind, and
it took the race to show it. One writer at a time, the probe sees only the file
check, which GitLab's also passes. Only under simultaneous writers does the
second refusal appear: git's ref lock on the push.

\textbf{Compaction is the exception on every host.} It commits a parentless
snapshot, which is not a fast-forward of anything, so on GitHub the reference
has to be forced, and the refs API takes no expected value for a forced update.
Forcing it unconditionally would discard, without a word, a sync another writer
landed after this device last looked. The engine therefore reads the reference
immediately before the forced move and refuses if it has moved. That turns a
silent loss into a refusal but is itself a check-then-act, so compaction is
safe for a machine with one writer at the time it runs and should be scheduled
as such. GitLab and Forgejo cannot make a parentless commit at all, and
compaction there is a branch recreated by hand.

\textbf{The protocol therefore guarantees no more than its host does.} Where the
service offers an atomic primitive, phase P7 is a compare-and-swap and no landed
sync is lost. Where it offers only a check-then-act, the design is safe for
writers that hand off in sequence and unsafe for writers that overlap in time.
A deployment on such a host should serialize its writers rather than rely on
the protocol. Nothing in the engine can supply an atomicity the API does not
have.

\section{Against a Conventional Object Store}
\label{sup:objectstore}

The design rules of Section~\Mref{sec:rules} of the main text are for services
that meter requests, so which do? Amazon S3 meters bytes and requests without a
ceiling of the kind we hit, does not address by content, and since 2024 offers a
compare-and-swap on an object~\cite{s3conditional,s3cas2024}. IPFS addresses by content,
but its mutable name is a signed record resolved through a hash table, not a
compare-and-swap, and persistence depends on pinning~\cite{ipfs}. A forge alone
supplies content addressing, a compare-and-swap reference and durability
together, from an account the user already holds.

Table~\ref{tab:objectstore} runs our commit shape on both. The object-store
column is the same engine's shape against a store with no tree and no commit.
It issues a PUT per new chunk under its content hash, a PUT for the sync's
manifest under a per-sync key so the history is kept, and a conditional PUT on a
pointer object. Both columns are counted by the same accounting on the same workload, so the
difference is the substrate and nothing else.

\textbf{What this is and is not.} Both columns run in process against counting
hosts, so the request and byte counts are exact: they follow from what the
protocol must issue, not from a network. The claim is about request shape,
which is what the design's cost model turns on, and the price below is
arithmetic over a published rate card rather than a bill we were charged.

\textbf{The same shape over real HTTP.} A model whose two sides we both wrote
shows less than it appears to, so we ran the object-store shape again,
unchanged, against a server on the loopback: \texttt{src/analysis/s3-server.js}.
It implements GET, PUT, DELETE and the two conditional PUT forms S3 gained in
2024, with the status codes S3 documents for them. Those are 412 for a
create-if-absent on a present key or for a stale ETag, and 404 for a
replace-if-unchanged against a key that holds nothing. The shape does not know
which host it has. Measurement is taken from the server's log, and
Table~\ref{tab:wire} gives it. Every figure matches the in-process model, and
the client's count of what it sent matches the server's count of what arrived:
\R{wire.served} requests, including the restore after each of the \R{wire.syncs} syncs. Both count
the pointer object's own bytes.

\begin{table}[!htb]
\caption{The object-store shape in process and over HTTP, on the \R{wire.syncs}-sync
session. The right column is read from the server's access log.}
\small
\centering
\begin{tabular}{@{}lrr@{}}
\toprule
 & in process & over HTTP \\
\midrule
\gen{tab-wire.tex}
\bottomrule
\end{tabular}
\label{tab:wire}
\end{table}

\textbf{And what it still is not.} The server is ours and it is on the
loopback, which is how the restic comparison of Section~\ref{sup:restic} is measured and
for the same reason: it makes the request path real without pretending the
network is. Nothing here measures S3's latency, its throughput, its durability
or its behavior under load, and none of those should be inferred from it. We
had no credentials for a commercial object store and did not run against one.
The harness takes \texttt{--url} and will run against any store speaking these
operations: the experiment a reader with an account can perform and we could
not.

\begin{table}[!htb]
\caption{The same commit shape on a forge and on an object store, 1,000 syncs
on a 256\,MB disk. Restore is requests / MB fetched.}
\small
\setlength{\tabcolsep}{4pt}
\centering
\begin{tabular}{@{}rrrrr@{}}
\toprule
 & \multicolumn{2}{c}{restore} & \multicolumn{2}{c}{cumulative upload req} \\
\cmidrule(lr){2-3}\cmidrule(lr){4-5}
sync & forge & object store & forge & object store \\
\midrule
\gen{tab-objectstore.tex}
\bottomrule
\end{tabular}
\label{tab:objectstore}
\end{table}

\textbf{The difference is a constant, not a shape.} One request per restore,
for the reference and commit a forge needs to name a state where a store needs
one pointer; two per sync, for the tree and the commit. Storage differs by two
percent, which is the tree and commit objects and the per-chunk digests our
manifest carries. Restore is flat in both after the live set stops growing, and
the bytes are identical. Everything the evaluation concludes about commit
shape, bounded restore and request economy therefore belongs to any store that
holds one object per request, and only the constant belongs to git.

What does not transfer is the part that is not a request count. A forge has a
ceiling and no price; a store has a price and no ceiling of that kind.
\texttt{price.mjs} prices the 1\,GB machine of Section~\ref{sup:histories} from its archived
replay at S3 Standard's published us-east-1 rates~\cite{s3pricing}: \$0.005 per
thousand PUTs, \$0.0004 per thousand GETs, \$0.023 per GB-month, and data out
free to 100\,GB a month. An object store spends two requests fewer per sync than
our shape, the exact difference above. The history's \R{price.uploads} upload requests are
therefore \R{price.puts} PUTs, \R{price.perSync} a sync. A weekly restore is \R{price.gets} GETs for \R{price.restoreMB}\,MB,
and our \R{price.storedMB}\,MB of storage bounds the store's from above. Synced once a day with
that
average sync, the machine costs \R{price.daily} a month: \R{price.dailyStorage} of storage, \R{price.dailyWrites} of
writes and nothing measurable for reads. Synced as hard as the captured
workload, all \R{price.syncs} syncs every day, it costs \R{price.hard}, of which \R{price.hardWrites} is writes.
The bill is storage for a machine at rest and requests for one in use, and the
second is the ceiling in another currency. Pricing a whole-machine commit every
day instead gives \R{price.whole}, but that is a compaction rather than a sync.

\section{Failure and Recovery}
\label{sup:faults}

Table~\ref{tab:faults} fails a sync in each of its phases and measures what the
failure cost. We take the machine to a settled state, write four chunks, and
make the host fail at one phase. The client then does the only thing it knows
how to do, which is sync again. Finally the disk is restored from the branch
name alone and compared with the device byte for byte. The row \emph{the reply}
fails after the last phase: the host moves the reference and the answer saying
so is lost, which the client cannot tell from a failed update. The host is the
in-process counting host, and the \R{faults.encryptedCases} network and server rows are run again
with encryption on.

``Spent'' is the requests the failed attempt made, ``left'' the objects it put
in the store, ``reused'' those of them the retry referenced rather than sending
again, and ``recovery'' the requests the next sync needed. In every row the
restored disk was identical and the reference never moved before its phase.

\begin{table*}[t]
\caption{A sync failed in every phase, on an \R{faults.diskMB}\,MB machine with four dirty
chunks. No row lost a write; no row orphaned an object.}
\small
\centering
\setlength{\tabcolsep}{1pt}
\begin{tabular}{@{}llrrrrrl@{}}
\toprule
fails in & as & spent & left & reused & orphaned & recovery & lost \\
\midrule
\gen{tab-faults.tex}
\bottomrule
\end{tabular}
\label{tab:faults}
\end{table*}

Four things follow. Nothing an interrupted sync uploaded is wasted, because
the retry recomputes the same content ids and the tree names them. Nothing is
orphaned for the same reason, so the garbage a failure leaves is zero rather
than merely collectable. Under encryption a retry would compute new ids, since
every encryption draws a fresh nonce, so the engine keeps a failed attempt's
ciphertexts and resends one whose plaintext has not changed. Without that, a
retry orphaned everything the attempt had uploaded. The reference is the commit
point: a failure anywhere before it leaves the branch exactly where it was.
After it, the state is already published, and the retry is refused, since its
parent is no longer the head. The engine records the manifest each unanswered
attempt would have published. When a refusal finds that manifest at the head, the
commit is its own; it adopts the commit rather than refusing its epoch as
another writer's. That costs \R{faults.recoveryLost} requests rather than \R{faults.recovery}, the refused attempt and the
read of the head. Elsewhere the recovery cost does not depend on where the
failure happened. That leaves something unclaimed: a failure at the reference
update has a complete commit sitting in the store, and retrying the reference
alone would cost one request rather than eight. The engine does not, because a client that remembered a commit id across
a crash would have to decide whether it was still the right parent, and the
eight requests are cheaper than that decision.

\section{Scaling Past the Measured Sizes}
\label{sup:scaling}

The evaluation reaches 1\,GB, where the hourly request budget starts to bind.
Rather than assert what 8\,GB would do, we state the model, check it against
both machines we did measure, and extrapolate with its error attached.

A restore reads the distinct objects the manifest names plus a fixed number of
requests to find them: two for the reference on GitHub, one for the tree, one
for the manifest. A sync writes the new objects plus a tree batch per thousand
live objects, a commit and a reference update. Neither term involves the
history. Against the two captured machines, four times apart in disk and eight
in live data, the model predicts restore and commit requests within
\textbf{\R{model.requestError}\%} and restore wall clock within \textbf{\R{model.timeError}\%}; on GitLab, whose
tree listing pages, the request prediction is exact and the wall clock is
within a tenth of a percent once its slower per-request rate is used.

\begin{table}[!htb]
\caption{The model past the measured sizes, at 75\% occupancy, 256\,KB chunks
and the \R{model.duplicateShare}\% duplicate share the 1\,GB machine had. ``Hourly budget'' is a
restore plus a commit against a token's 5,000 requests an
hour~\cite{ghratelimits}.}
\small
\setlength{\tabcolsep}{4pt}
\centering
\begin{tabular}{@{}lrrrrr@{}}
\toprule
live disk & objects & restore & \hd{at 8\\ways} & commit & \hd{hourly\\budget} \\
\midrule
\gen{tab-scalingmodel.tex}
\bottomrule
\end{tabular}
\label{tab:scalingmodel}
\end{table}

\textbf{The wall is requests, not bytes or storage.} At 256\,KB chunks one
hour of a token's allowance restores about \R{model.oneHourGB}\,GB of live disk. A 6\,GB
machine needs about \R{model.sixRestoreWindows} hourly windows for a restore and \R{model.sixBothWindows} for a restore and a
commit together, while its restore's bytes take \R{model.sixWideMinutes} minutes at eight ways. Storage is not
the constraint either: 6\,GB of live disk is inside GitLab's 10\,GB per project
and past what GitHub recommends, which is a policy limit rather than a
mechanical one.

Two levers move the wall, and the table prices both. Doubling the chunk halves
every request count, so the same 6\,GB machine restores in \R{model.six512} requests at
512\,KB and \R{model.six1M} at 1\,MB, at the byte cost Section~\ref{sup:chunk} measures. And a
batch-commit host removes the commit side entirely, since its sync is one
request whatever the machine's size, which is the portability result of
Section~\Mref{sec:portability} of the main text turning from an economy into a scaling property.




\section{Deterministic Encryption}
\label{sup:deterministic}

This section explains why the engine encrypts under random nonces rather than
deterministically. A deterministic variant derives each chunk's nonce from its
plaintext under the key, so identical chunks encrypt identically and free space
still collapses to one object. That is deterministic authenticated encryption in
the sense of Rogaway and Shrimpton~\cite{siv}: it leaks plaintext equality and
nothing else. On a chunk-aligned block store, equality has three classes. One
is free space, the reason for the variant. Another is a chunk reverting to
earlier content, which reveals a rollback of that chunk. The third is duplicates
across positions, which Section~\Mref{sec:dedup} of the main text found only in a
format's structure and in copies a user makes.

The hazard is the confirmation attack of Harnik et al.~\cite{harnik}, which
needs no key. An adversary who can read the repository and get content onto the
disk, say a file the user downloads, learns by ciphertext equality alone whether
it was already there. Alignment raises the bar, since the bytes must fill a
256\,KB chunk exactly and the observer must be a reader who is not the owner, as
a public or shared repository supplies. Yu et al.~\cite{sidechannel} treat the
multi-tenant form, and the exchange that followed~\cite{sidechannelcomment}
shows its defences unsettled.

Random nonces remove the hazard rather than narrowing it, and free space keeps
its saving: the engine represents a chunk the guest has zeroed on a blank base
by absence, exactly as an unwritten chunk. It encrypts every stored chunk with
AES-GCM~\cite{gcm} under a fresh random 96-bit nonce, with the chunk index and salt as
associated data, so an object decrypts only at its own position. The key is
derived from a passphrase by PBKDF2-HMAC-SHA256~\cite{rfc8018} over 250,000 iterations. No two
ciphertexts can be related, and a host cannot rearrange them.

The derivation yields two 256-bit blocks. The first is the chunk key, and the
second keys an HMAC-SHA256 over every manifest the machine commits, all of it
but the MAC itself. Attach and restore verify it before the manifest is used,
so a host without the passphrase can change nothing in the manifest: not which
object sits at which index, not a digest, not the format version. What remains
is a whole older manifest, which the pinned digest refuses. A cipher that does
not fit the machine is refused before anything is written: no passphrase for an
encrypted machine, one derived for another salt, or one given for a plaintext
machine. The residual leak,
occupancy and write pattern, is what the commit tree reveals under any cipher.

What random nonces give up is every other duplicate. Table~\ref{tab:encrypted}
replays both captured machines through the real engine twice, in plaintext and
encrypted. The working session repeats only its format's chunks and a package
it removes and adds back, and loses little. The corpus workload copies files
with \texttt{tar}, and loses a quarter. The cipher itself is cheap: in
\R{enc.runtime}, a key takes \R{enc.deriveMs}\,ms to derive and a chunk \R{enc.chunkMs}\,ms to encrypt, medians in one
process.

\begin{table}[!htb]
\caption{Encryption's cost at the last sync of each captured machine:
requests to restore, upload requests over the whole history, and MB stored, in
plaintext and under random nonces, on the in-process host. Both restored byte
for byte. The 1\,GB row is a second capture of the corpus workload, kept with
its bytes. Its live set equals the first capture's, with \R{census.big.distinct} distinct objects
against the first's \R{census.big.originalDistinct}.}
\small
\setlength{\tabcolsep}{3pt}
\centering
\begin{tabular}{@{}lrrrrrr@{}}
\toprule
& \multicolumn{2}{c}{restore} & \multicolumn{2}{c}{uploads} & \multicolumn{2}{c}{MB stored} \\
\cmidrule(lr){2-3}\cmidrule(lr){4-5}\cmidrule(lr){6-7}
machine & plain & enc. & plain & enc. & plain & enc. \\
\midrule
\gen{tab-encrypted.tex}
\bottomrule
\end{tabular}
\label{tab:encrypted}
\end{table}

The general lesson is narrow. Deterministic encryption is defensible only with
one key per tenant, no adversary able to both inject plaintext and read
ciphertext, and equality classes of no value. Deduplication across tenants under
a shared key fails all three. Table~\ref{tab:threat} summarizes what is
protected, against whom, and what remains.

\section{restic Repeats}
\label{sup:restic}

The measured tools are restic 0.19.1~\cite{resticsw} and rest-server
0.14.0~\cite{restserver}. restic's chunker is seeded per repository. We repeated the 120-sync run on a
64\,MB disk, the default of the harness, three times. Restore at sync 120 took
\R{repeat.req} requests for \R{repeat.mb}\,MB. Storage was \R{repeat.storage}\,MB,
and each run made \R{repeat.uploads} upload requests. The runs agree within \R{repeat.reqSpread}
request, $\pm$\R{repeat.bytesSpread}\% of bytes and $\pm$\R{repeat.storageSpread}\% of storage about the middle of each
range.

Restore is by a client with no cache. The \texttt{repair index} figures in
Section~\Mref{sec:scaling} of the main text are the 1,000-sync run on the
256\,MB disk. The rewrite itself costs \R{restic.repaired.rewrite} requests. Restore then reads \R{restic.repaired.ranges}
ranges to cover a live set of \R{scale.ours.live} chunks and costs \R{restic.repaired.req} requests for \R{restic.repaired.mb}\,MB,
against our \R{scale.ours.req} for \R{scale.ours.mb}\,MB.

On a forge every one of those ranges would be a whole pack, since its object API
has no ranged read for a private repository. We reran restic on both captured
histories and recorded, for each restore, the packs it read. Compacted, it read
\R{rf.hist.packs} packs at sync 43: \R{rf.hist.rangeMB}\,MB by range, or \R{rf.hist.wholeMB}\,MB in \R{rf.hist.wholeReq} requests read whole. At sync
\R{long.syncs} of the long session it read \R{rf.long.packs} packs: \R{rf.long.rangeMB}\,MB by range, or \R{rf.long.wholeMB}\,MB in
\R{rf.long.wholeReq} requests. Its storage at that point is \R{restic.repaired.storageMB}\,MB against
our \R{scale.ours.storageMB}\,MB.

% The supplement cites the same database as the main text. Its own
% reference list is short and numbered independently, so a citation here
% never has to agree with a number over there.
{\footnotesize
\bibliographystyle{IEEEtran}
\bibliography{refs}
}

\end{document}
